\section{Introduction}
\label{sec:introduction}
Language agents use external tools to carry out multi-step tasks \citep{react,schick2023toolformer}. Benchmarks cover shopping, travel, and service workflows \citep{webshop,travelplanner,taubench}, as well as broader web tasks \citep{zhou2024webarena,deng2023mind2web}. ST-WebAgentBench counts completion as successful only when applicable policies are satisfied \citep{stwebagentbench}. In delegated purchasing, the commitment and its possible consequences must stay within the user's authorization. Cancellation, return, or replacement costs can exceed that authorization even when the initial price is within budget.

Requests can have several readings, and future events determine financial consequences. We study purchase or refusal without another user turn, using only existing authorization. Choosing a convenient reading can permit an unauthorized purchase; refusing every uncertain request can discard useful options. Yet one purchase may satisfy every retained possibility. The task is to find a common authorized purchase with positive worst-case value.

Clarification and uncertainty-aware planning use additional information to improve decisions \citep{knowno,voi,sage}. We apply the fixed-set max--min principle of Risk-Averse Calibration (RAC) \citep{rac}, with purchasing permissions that cannot be traded against utility. Turning intent into executable decisions requires preserving which conditions apply under each reading or event. For example, a loss limit restricts permission; it is not compensation that reduces the loss. That distinction affects both admissibility and value.

We introduce \ours\ to evaluate authorization across retained states and rank admissible purchases by worst-case utility through a fixed procedure. We assess how reliably this workflow turns natural-language delegation into purchasing decisions across model families.

A preferred offer becoming unavailable does not expand the user's authorization. Yet a safe purchase can reflect either a rule that excludes unauthorized offers or a ranking that keeps them from winning. In one controlled replay, a specification initially makes an exact decision, but selects an unauthorized offer after four high-ranked offers are removed, while 59 alternatives remain permitted and positively valued under both its conditions and the reference. The request, permissions, utilities, and surviving offers are unchanged. We analyze how availability and valuation separately expose existing permission errors, and why a correct current purchase cannot certify the entire admitted domain.

The paper makes three contributions:
\begin{itemize}
\item \textbf{Risk-aware decisions under unresolved uncertainty:} a formulation of purchase or refusal that requires a common authorized action and separates its permission from its worst-case value.
\item \textbf{Authorization errors hidden by selection:} we distinguish selected-action correctness from domain soundness, derive the effects of state omission, and explain when availability or valuation exposes an existing permission error.
\item \textbf{Controlled empirical investigation:} multi-model boundary and intent tests assess decision reliability; a prospective availability study retains viable authorized alternatives, and separate component interventions trace permission--value interactions.
\end{itemize}

Code, benchmark materials, retained responses, evaluation records, and reproduction tools are available in our anonymous artifact package (Appendix~\ref{app:reproduction}).\footnote{\url{https://anonymous.4open.science/r/risk-aware-payment-delegation-D4FC/}}
